% ABOUTME: Introduction section of the IEEE report.
% ABOUTME: Motivation for RUL prediction, the C-MAPSS benchmark, our approach, contributions and paper layout.
\section{Introduction}
\label{sec:intro}

Maintenance of aircraft engines is either scheduled at fixed intervals or
triggered by observed condition. Fixed intervals replace parts that still
have life left, while waiting for a fault risks a failure in service.
Predicting the remaining useful life (RUL) of an engine, the number of
operating cycles left before it fails, lets maintenance be planned close to
the point where it is actually needed. The cost of an error is asymmetric:
predicting too much remaining life risks an in-service failure, while
predicting too little only brings maintenance forward.

The NASA Commercial Modular Aero-Propulsion System Simulation (C-MAPSS)
data~\cite{saxena2008damage} are the standard public benchmark for this
task. Each engine is simulated from a healthy state with unknown initial
wear, and 21 sensors and three operating settings are recorded once per
cycle. Training engines run to failure; test engines are cut off at an
unknown point, and the true RUL is given only at their last observed cycle.
We work on the FD001 subset, which has one operating condition and one
fault mode.

Our approach keeps the learners simple and puts the effort into the
features and the protocol. From the 14 sensors that vary in FD001 we
compute, within each engine, rolling means, standard deviations and least
squares slopes over three window lengths, and a health index learned by
linear regression. Elastic Net, Random Forest and XGBoost are trained on
RUL capped at 125 cycles and tuned with five-fold cross validation grouped
by engine, so that no engine contributes rows to both sides of a split. The
test set is used once, after every hyperparameter and the choice of model
are frozen.

The contributions of this report are:
\begin{itemize}
\item an engine-level evaluation protocol for FD001 in which feature
statistics are refitted inside every fold and the test set is scored once,
with errors reported separately for engines near failure;
\item a comparison of three regressors on raw and engineered features,
showing that the engineered features improve every model and account for
most of the gain over raw sensors;
\item an analysis of what the tree ensembles rely on, using held-out
grouped permutation importance and SHAP values, and a comparison of the
test results with published FD001 models under matched label clipping.
\end{itemize}

Section~\ref{sec:related} reviews prior work on C-MAPSS.
Section~\ref{sec:data} describes the data and the choice of sensors and
target cap, and Section~\ref{sec:evaluation} the evaluation protocol.
Section~\ref{sec:features} defines the engineered features.
Section~\ref{sec:models} covers the models, their tuning and what they rely
on, and Section~\ref{sec:results} reports the test results.
Section~\ref{sec:discussion} discusses the findings and their limits, and
Section~\ref{sec:conclusion} concludes with the plan for sequence models.
